\chapter{Related Work} \label{ch:related}

This chapter positions our work relative to prior research on Bitcoin vault custody,
covenant design, formal verification of cryptocurrency protocols, and layer-two
security.

% ======================================================================
\section{Vault Custody Protocols} \label{se:rw_vault_custody}

Swambo et al.~\cite{SHMB20} formalized the vault lifecycle model---deposit, unvault,
withdraw, recover---and the watchtower monitoring assumption that underpins all
covenant vault designs. Their analysis established the threat vocabulary (key
compromise, liveness denial, fund safety) that we adopt and extend with the
11-threat taxonomy (TM1--TM11). A companion paper~\cite{SBMB20} characterized three
approaches to implementing covenants (pre-signed transactions with deleted keys,
OP\_CHECKOUTPUTSHASHVERIFY, and OP\_CHECKTEMPLATEVERIFY), analyzing the security
assumptions of each. Their work analyzed deleted-key covenants and pre-signed
transaction vaults without comparing opcode-level designs empirically. Swambo's
doctoral thesis~\cite{Swa23} extended this to the Ajolote autonomous custody system
with formal security analysis, but did not include CTV, CCV, OP\_VAULT, or
CAT+CSFS.

M\"oser, Eyal, and G\"un Sirer~\cite{MES16} introduced covenants as a mechanism for
vault-based custody, proposing script-level restrictions on how UTXOs can be spent.
Their work predates all four BIP proposals studied here and provides the conceptual
foundation for the covenant approach to custody. O'Connor and
Piekarska~\cite{OP17} demonstrated that covenants can be implemented in Bitcoin by
adding OP\_CAT and OP\_CHECKSIGFROMSTACKVERIFY, foreshadowing the CAT+CSFS design
studied in this thesis.

O'Beirne's technical report on vaults and covenants~\cite{OBeirne23} provides a
practitioner-oriented comparison of vault designs, analyzing the tradeoffs between
OP\_VAULT's purpose-built approach and CTV's generic template commitment. This
informed the design of BIP-345~\cite{BIP345}, which was subsequently
withdrawn~\cite{OB25} in favor of CCV (BIP-443).

Kirstein et al.~\cite{KLGZ21} proposed Phoenix, a formally verified regenerating
vault that recovers from key compromise by regenerating the vault with new keys.
Their work uses formal verification (in Dafny) to prove safety properties of a
single vault design. Our work differs by comparing five designs under a uniform
methodology rather than deeply verifying one.

The Bitcoin PIPEs v2 proposal~\cite{PIPES26} introduces an alternative approach to
covenants using zero-knowledge proofs, enabling programmable vaults without new
opcodes. This represents a fundamentally different design philosophy from the
opcode-level covenants studied here.

% ======================================================================
\section{Covenant Proposals and Analysis} \label{se:rw_covenants}

Each covenant proposal has its own design documentation and community analysis, but
no prior work compares them empirically.

Rubin~\cite{BIP119} specified CTV and discussed its applications (vaults, congestion
control, payment pools) but did not provide measured transaction sizes or cross-design
comparisons. The bitcoin-dev mailing list discussion identified address reuse as a
risk, which we confirm empirically.

Ingala~\cite{BIP443} designed CCV (originally MATT) and documented the keyless recovery
griefing concern and the OP\_SUCCESS behavior for undefined modes. Our work quantifies
the griefing cost asymmetry ($\gamma = 0.27$) and confirms the OP\_SUCCESS behavior
across five undefined mode values on a production-shaped taptree.

O'Beirne and Sanders~\cite{BIP345} specified OP\_VAULT with authorized recovery and
automatic revault. Harding~\cite{Har24} estimated per-UTXO recovery-cost scaling at
approximately 3{,}000 splits per block. We confirm this for OP\_VAULT (3{,}427
measured) and extend it to CCV (6{,}172 splits per block, 80\% more than Harding's
estimate). O'Beirne subsequently withdrew BIP-345~\cite{OB25}; our lifecycle cost
comparison (OP\_VAULT 36\% more expensive than CCV) provides quantitative context for
that decision.

Poelstra~\cite{Poe21} demonstrated transaction introspection via OP\_CAT and Schnorr
signature tricks, constructing a vault with recursive covenants and dynamic destination
selection. Our CAT+CSFS vault implements a simplified variant (fixed destination, no
recursion) to isolate the dual-verification security properties. The cold key recovery
vulnerability (TM11) and the Poelstra reset cost model in Chapter~\ref{ch:security}
directly build on this work. Rijndael~\cite{Rij24} built a working vault prototype
using only OP\_CAT and the Schnorr G-trick (without CSFS), demonstrating that OP\_CAT
alone enables covenant enforcement.

% ======================================================================
\section{Formal Verification of Bitcoin Contracts} \label{se:rw_formal}

Bartoletti and Zunino~\cite{BZ18} introduced BitML, a domain-specific calculus for
Bitcoin smart contracts with formal semantics based on process algebra. BitML enables
reasoning about contract security through symbolic and computational models.
Subsequent work developed a toolchain for compiling BitML to Bitcoin
transactions~\cite{ABL19} and verified liquidity properties of recursive
contracts~\cite{BLM22}. Bartoletti, Lande, and Zunino~\cite{BLZ20} extended this
line to Bitcoin covenants, modeling vault contracts in an extended BitML with
covenant primitives and proving safety properties. Our Alloy models operate at a
different abstraction level---state-machine reachability rather than process-algebraic
semantics---but address the same class of questions: can an attacker extract funds
from a correctly constructed vault?

dgpv~\cite{dgpv24} applied Alloy bounded model checking to Rijndael's OP\_CAT vault
prototype, verifying fund conservation and covenant enforcement properties within
a single design. Our work extends this methodology to all four BIP-specified
designs plus Simplicity, using a shared base model (\texttt{vault\_base.als}) to
ensure uniform property definitions. We add the threat model layer
(\texttt{threat\_model.als}) for systematic attacker profiling and the cross-covenant
module (\texttt{cross\_covenant.als}) for multi-design interaction
checking---neither of which appears in dgpv's analysis.

O'Connor~\cite{OC17} provided Coq-verified denotational semantics for the Simplicity
language, establishing the formal foundation for reasoning about Simplicity programs.
Our Alloy model of the Simplicity vault operates at the state-machine level (vault
lifecycle transitions) rather than the language semantics level, complementing
O'Connor's foundational work with application-level verification. O'Connor and
Poelstra have also pursued formal verification of cryptographic primitives used in
Bitcoin's libsecp256k1 library using Coq, demonstrating the viability of machine-checked
proofs for Bitcoin infrastructure.

Several surveys have cataloged formal verification methods for blockchain smart
contracts. Murray and Anisi~\cite{Murray19} categorized approaches including model
checking, theorem proving, and symbolic execution. Tolmach et al.~\cite{Tolmach21}
provided a comprehensive survey of formal specification and verification techniques,
covering 50+ tools and frameworks. Almakhour et al.~\cite{Almakhour20} focused on
verification methods applicable to production smart contracts, and Singh
et al.~\cite{Singh20} analyzed how formalization approaches address known
vulnerability classes. Our use of Alloy bounded model checking falls within the model
checking category identified by these surveys, applied specifically to Bitcoin
covenant protocols rather than Ethereum smart contracts.

% ======================================================================
\section{Layer-Two Security and Watchtower Economics} \label{se:rw_layer2}

The watchtower monitoring model used in vault security is closely related to the
watchtower concept in Lightning Network payment channels. Leinweber
et al.~\cite{LGSM19} proposed TEE-based distributed watchtowers for fraud protection
in Lightning, analyzing the economic viability of third-party watchtower services.
Our per-UTXO recovery-cost scaling analysis (TM4) addresses a different question---how an
attacker can exhaust the watchtower's recovery budget through vault
splitting---but shares the assumption of an always-online monitoring service.

Naumenko and Riard~\cite{NR21} analyzed time-dilation attacks on the Lightning
Network, where an attacker with network-level capabilities manipulates the victim's
view of the blockchain to compress timelocks. This attack class applies to covenant
vaults (an attacker could compress the CSV window), but we do not model it because
it is orthogonal to covenant-specific properties and requires network-layer
simulation beyond our regtest methodology. Riard~\cite{Ria23} subsequently
demonstrated replacement cycling attacks against Lightning HTLCs, a mempool
manipulation technique that may also apply to covenant vault recovery transactions.

Aumayr~\cite{Aumayr24} provided a comprehensive formal treatment of
Bitcoin-compatible scalability protocols, including watchtower collateralization
and channel closing security. The economic analysis of watchtower reserves in that
work parallels our per-UTXO recovery-cost scaling budget calculations, though applied to
payment channels rather than vaults.

% ======================================================================
\section{Qualitative Comparisons} \label{se:rw_qualitative}

The covenants.info website~\cite{covenants_info} provides a qualitative feature matrix
comparing covenant proposals across dimensions such as recursion support, introspection
granularity, and application breadth. This resource contains no transaction sizes, no
attack cost functions, and no fee-environment analysis. Our work complements it with
quantitative measurements.

The B-SSL covenant-free vault model~\cite{BSSL25} demonstrates vault-like custody
using only CSV and CLTV timelocks (no new opcodes). This provides a baseline for
what is achievable without covenants---the security properties are weaker (no output
constraints on recovery, no dynamic destination selection) but the deployment path
requires no consensus changes.

% ======================================================================
\section{Summary of Positioning} \label{se:rw_positioning}

Table~\ref{tab:related_work} compares this work with the most closely related prior
efforts across five dimensions.

\begin{table}[t]
\centering
\caption{Comparison with closely related prior work.}
\label{tab:related_work}
\small
\begin{tabular}{@{}lccccc@{}}
\toprule
Work & Covenants & Empirical? & Formal? & Fee analysis? & Year \\
\midrule
M\"oser et al.~\cite{MES16}      & 1 (generic)
  & No  & No & No  & 2016 \\
Swambo et al.~\cite{SHMB20}  & 0\textsuperscript{a}
  & No  & Partial & No  & 2020 \\
Bartoletti et al.~\cite{BLZ20}   & 1 (BitML)
  & No  & Yes (process) & No  & 2020 \\
Kirstein et al.~\cite{KLGZ21}    & 1 (Phoenix)
  & No  & Yes (Dafny) & No  & 2021 \\
Harding~\cite{Har24}          & 1 (OPV)
  & No\textsuperscript{b}  & No & Partial & 2024 \\
dgpv~\cite{dgpv24}            & 1 (CAT)
  & No  & Yes (Alloy) & No  & 2024 \\
covenants.info~\cite{covenants_info} & 6+
  & No  & No & No  & 2024 \\
\textbf{This work}            & \textbf{5}
  & \textbf{Yes} & \textbf{Yes (Alloy)} & \textbf{Yes} & \textbf{2026} \\
\bottomrule
\multicolumn{6}{@{}l}{\footnotesize \textsuperscript{a}Analyzes vault model
generically, not specific BIP opcodes.} \\
\multicolumn{6}{@{}l}{\footnotesize \textsuperscript{b}Estimates from assumed
transaction sizes, not regtest measurements.}
\end{tabular}
\end{table}

No prior work combines empirical measurement, formal verification, and
fee-environment-dependent security analysis for multiple covenant vault designs.
The individual components---vault lifecycle modeling~\cite{SHMB20}, Bitcoin contract
formalization~\cite{BZ18, BLZ20}, Alloy verification~\cite{dgpv24}, watchtower
exhaustion estimation~\cite{Har24}---exist in isolation. Our contribution is the
integration of these approaches under a uniform methodology that enables cross-design
comparison and the derivation of results (Proposition~\ref{thm:griefing_safety}
and~\ref{thm:fee_inversion}) that require data from multiple designs simultaneously.
